\section{Detalle completo del juego de integración}
\label{app:megajuego-full-detail}

La Sección~\ref{subsec:megajuego} enuncia cuatro resultados en formato
compacto; sus demostraciones completas siguen a continuación.

\subsection{Potencial exacto y óptimo central}
\label{app:megajuego-potential-proof}

\begin{proof}[Demostración de la Proposición~\ref{prop:megajuego-exact-potential}]
Los factores no incidentes no cambian y se cancelan. No se define el coste
social como $\sum_iJ_i$: eso duplicaría factores compartidos. Si se eliminan
los estados mediante integración, una entrada local puede afectar costes
remotos; esos gradientes deben propagarse mediante adjuntos o mensajes, o
acotarse. La factorización antes de eliminar estados no prueba localidad de
la sensibilidad reducida.
\end{proof}

\begin{proof}[Demostración del Teorema~\ref{thm:megajuego-central-optimum}]
Es la condición de primer orden de minimización convexa; en ambos sentidos se
utiliza el segmento factible entre $z^\star$ y $z$.
\end{proof}

Un GNE ordinario no basta: con $z_1+z_2=1$, cada jugador puede quedar sin
desviaciones unilaterales y todo punto factible ser GNE. Se busca la
desigualdad variacional (VI) o revisiones conjuntas suficientes. Fijar
coalición y homotopía no garantiza por sí solo la convexidad requerida.

\subsection{Separación exacta de trayectorias Bézier cooperativas}
\label{app:megajuego-bezier-proof}

Se consideran dos coaliciones Cargo, cada una contenida en un disco exterior
de radio $1.05$~m; la separación adicional es $0.1$~m, de modo que
$D=2.2$~m. En $s=t/20\in[0,1]$, $P_A(s)=(-3+6s,\,a_Ab(s))$,
$P_B(s)=(3-6s,\,-a_Bb(s))$, $b(s)=64s^3(1-s)^3$. No hay otros obstáculos; el
horizonte y el avance longitudinal están fijados.

\begin{proof}[Demostración del Teorema~\ref{thm:megajuego-bezier-separation}]
En $s=1/2$, la distancia es $a_A+a_B$, lo que prueba necesidad. Sea
$z=(2s-1)^2\in[0,1]$ y $A=a_A+a_B\geq D$. Entonces
\[
 \|P_A-P_B\|^2=36z+A^2(1-z)^6\geq36z+D^2(1-z)^6.
\]
La derivada del último miembro respecto a $z$ es $36-6D^2(1-z)^5\geq0$ si
$D\leq\sqrt6$. Su mínimo es $D^2$ en $z=0$.
\end{proof}

Este certificado es para todo tiempo, no una ausencia de colisión muestreada.
$b$ es una curva Bézier de grado seis con coeficientes
$(0,0,0,16/5,0,0,0)$; además $b'(0)=b'(1)=b''(0)=b''(1)=0$, de modo que todas
las amplitudes comparadas comparten poses, velocidades longitudinales y
velocidades iniciales de ruedas.

El objetivo es la suavidad transversal ponderada
$\Phi(a)=\tfrac12h_Aa_A^2+\tfrac12h_Ba_B^2$ sujeto a $a_A+a_B\geq D$,
$a_i\geq0$, relacionado exactamente con la aceleración transversal cuadrática
porque $\int_0^{20}\ddot y_i^2\,dt=\tfrac{8192}{35\cdot20^3}a_i^2$. Los pesos
positivos producen óptimo en $a_A+a_B=D$, con
$a_A^\star=h_BD/(h_A+h_B)$, $a_B^\star=h_AD/(h_A+h_B)$,
$\Phi^\star=h_Ah_BD^2/(2(h_A+h_B))$. Para $(h_A,h_B)=(2,1)$,
$a_A^\star=11/15$, $a_B^\star=22/15$; la solución simétrica
$a_A=a_B=1.1$ cuesta $363/200$, un $12.5$\,\% más. Si cada jugador minimiza
solo su coste privado sujeto a $a_A+a_B\geq D$, cualquier reparto en la
frontera es un GNE ordinario: nadie puede reducir unilateralmente su
amplitud. El vGNE añade el multiplicador común y exige $h_Aa_A=h_Ba_B$,
seleccionando precisamente el reparto óptimo.

\paragraph{Protocolo distribuido ejecutado}
Cada coalición conoce su $h_i$, su propuesta y los mensajes de la otra. El
coordinador de incidencia usa $B=(1,-1)^{\mathsf T}$, $L=BB^{\mathsf T}$,
$b_i=D/2$ y estado $w=(a_A,a_B,\nu_A,\nu_B,\eta)$, con operador afín
\[
 \mathcal F(w)=
 \begin{bmatrix}\operatorname{diag}(h) & I & 0\\ -I & L & B\\ 0 & -B^{\mathsf T} & 0\end{bmatrix}w
 +\begin{bmatrix}0\\0\\D/2\\D/2\\0\end{bmatrix}.
\]
Se aplica extragradiente con paso $0.8/\|\mathcal F_{\mathrm{linear}}\|_2$ y
proyección local $a_i\in[0,D]$. La parte simétrica del operador es positiva
semidefinida y su solución primal es el óptimo anterior. En el caso principal
se alcanzó residual menor que $10^{-9}$ en 113 iteraciones; en 30 instancias
de pesos heterogéneos, todas alcanzaron esa tolerancia, con máximo
$9.998\times10^{-10}$ e intervalo 109--315 iteraciones. El óptimo analítico
se usa solo para evaluar el error; no interviene en las actualizaciones ni en
la condición de parada, que comprueba residual natural KKT, balance y
consenso, y brecha primal--dual de la rama.

\paragraph{Levantamiento físico ideal}
Cada carga tiene 40~kg y cuatro pads simétricos en $(\pm0.45,\pm0.45)$~m; el
yaw de carga es constante y los chasis se alinean con su tangente. Los robots
se modelan contenidos en discos exteriores de radio 0.32~m; la carga en un
rectángulo de $1.4\times1.3$~m; la unión queda dentro de la envolvente de
1.05~m. Con rodadura ideal, pads rígidos centrados y yaw libre, y torque como
entrada, el reparto $f_i=(m_\ell/4)\ddot P$ satisface Newton--Euler de carga
y momento cero por simetría. Las cotas de fuerza, ruedas y energía obtenidas
para $a\leq22/15$ son conservadoras sobre todo el intervalo; el máximo entre
los ocho AMR fue: velocidad de rueda $4.80260$~rad/s, torque por rueda
$0.325989$~N\,m, fuerza horizontal por pad $0.880001$~N, fuerza por rueda al
suelo $3.29597$~N y energía de batería por AMR en 20~s, $181.543$~J. Estas
cotas cubren geometría, ruedas, reacciones, torque y batería; no son una
simulación de \emph{tracking} de pads deformables ni de inductancia de motor.

\begin{proof}[Demostración del Corolario, óptimo físico nominal dentro de la familia]
Restringir el conjunto no puede disminuir su mínimo. Como el mínimo global
$a^\star$ de $\Phi$ en la familia geométrica posee el levantamiento factible
construido, $a^\star$ pertenece también al subconjunto que admite ruedas,
reacciones, torque y batería de la reducción ideal, luego ese mismo valor
sigue siendo alcanzable y ambas cotas coinciden. Este razonamiento respalda
optimalidad física nominal para el coste de suavidad y la familia fijados; no
sustituye ese coste por energía de batería ni extiende el resultado a otra
homotopía, duración, adquisición o planta con \emph{compliance}.
\end{proof}

\subsection{Caging: certificado de contacto con KKT racional}
\label{app:megajuego-caging-proof}

Se considera una carga circular de radio $R=0.45$~m y cuatro contactos
cardinales; el cerco tiene radio $0.65$~m y el robot un radio ocupado de
0.20~m. Con $z=(n_1,t_1,\ldots,n_4,t_4)$, $0\leq n_i\leq30$~N,
$|t_i|\leq0.4n_i$, se demanda un \emph{wrench} $(20,0,5)$ en N, N y N\,m; con
empujes solo normales y radiales el momento es siempre cero y esa demanda es
imposible. Con adherencia tangencial de Coulomb se resuelve
\[
 \min_z\tfrac12z^{\mathsf T}z \quad\text{sujeto a}\quad Az=(20,0,5)^{\mathsf T},\ Hz\leq g,
\]
con
\[
 A=\begin{bmatrix}-1&0&0&1&1&0&0&-1\\0&-1&-1&0&0&1&1&0\\0&-9/20&0&-9/20&0&-9/20&0&-9/20\end{bmatrix}.
\]
La solución racional es
\[
 z^\star=\Big(0,\,0,\,\tfrac{475}{261},\,-\tfrac{190}{29},\,\tfrac{500}{29},\,-\tfrac{200}{29},\,\tfrac{2275}{261},\,-\tfrac{910}{261}\Big),
\]
de valor $513025/2349$. El código verifica exactamente factibilidad,
estacionariedad y multiplicadores de desigualdad no negativos con SymPy. Por
convexidad estricta, es el óptimo único de este QP. La solución numérica
independiente se conserva con sus residuos, incluidos signos de brecha de
tamaño $10^{-11}$ debidos al redondeo.

\paragraph{Relación exacta con energía eléctrica estática}
Para robots estacionarios orientados hacia sus normales, radio de rueda
0.1~m, semivía y brazo de \emph{bumper} 0.2~m,
$\tau_{iL/R}=0.05(n_i\mp t_i)$. Con $nk_t=0.5$~N\,m/A y resistencia
$0.6\,\Omega$, $P_{cu}=0.6\sum_{i,s}(\tau_{is}/0.5)^2=0.024(z^{\mathsf T}z/2)$.
El mismo QP minimiza por tanto \emph{exactamente} las pérdidas de cobre en esa
rama estática; el óptimo es $P_{cu}^\star=20521/3915=5.241634738$~W. El
torque máximo es $1.206897$~N\,m, la corriente máxima $2.413794$~A y la
tensión estática máxima $1.448276$~V; la fuerza máxima por rueda requerida es
$12.551914$~N, inferior a $61.3125$~N disponibles. Un \emph{wrench} externo
opuesto mantiene el equilibrio estático de la carga; este experimento no es
transporte, demuestra un bloque físico--económico exacto y su certificado,
y muestra que un robot con fuerza nula puede seguir formando parte del cerco.

\subsection{Presupuesto de ejecución}
\label{app:megajuego-budget-proof}

\textbf{Definiciones usadas en la prueba.} Un \emph{tubo} es un intervalo
maximal de tiempo durante el cual la continuación vigente no cambia; una
ejecución es la concatenación de tubos $[t_0,t_1),[t_1,t_2),\ldots$, separados
por revisiones. Un \emph{testigo inicial seguro} es un certificado de que el
estado en $t_0$ admite al menos una continuación admisible (en el sentido de
la Proposición~\ref{prop:megajuego-exact-potential}: dinámica, contactos,
soporte, ruedas, energía, seguridad y causalidad respetadas). Que ``los tubos
componen'' significa que si el estado al final de un tubo porta un testigo
seguro y la revisión que le sigue es admitida, el estado al inicio del
siguiente tubo también porta un testigo seguro; esta es una hipótesis del
teorema, no una consecuencia que se derive de otro resultado del documento.

\begin{proof}[Teorema~\ref{thm:megajuego-execution-budget}]
Sea $T_{\mathrm{mov}}$ el tiempo acumulado de movimiento, es decir la medida
del conjunto de instantes en que la hipótesis de decrecimiento mínimo está
activa. Los intervalos sin movimiento no entran en $T_{\mathrm{mov}}$ y no
consumen $\kappa_W$; esto es una restricción del enunciado, no un tecnicismo:
una ejecución que espere indefinidamente sin moverse no queda acotada por este
resultado.

\emph{Cota de revisiones.} Cada sustitución admitida reduce el presupuesto en
al menos $\delta>0$, y el presupuesto es no negativo y parte de $W_0$, luego
$\delta N_{\mathrm{rev}}\leq W_0$ y $N_{\mathrm{rev}}\leq\lfloor W_0/\delta\rfloor$.
Esta cota no requiere suponer ausencia de Zeno: el número de
revisiones es finito por agotamiento del presupuesto, y de ahí que la sucesión
de tubos sea finita. Suponer además un margen mínimo entre revisiones sería
incorrecto, porque la ausencia de Zeno no implica un tiempo de permanencia
uniformemente positivo: una sucesión con $t_{k+1}-t_k=1/k$ no acumula
infinitos eventos en tiempo finito y aun así tiene margen ínfimo nulo.

\emph{Cota de tiempo de movimiento.} Por inducción sobre los
$k=0,\ldots,N_{\mathrm{rev}}$ tubos, que ya sabemos finitos. \emph{Base:}
$k=0$ porta un testigo seguro por hipótesis y el presupuesto inicial es $W_0$.
\emph{Paso:} si el tubo $k$ empieza con testigo seguro y presupuesto
$W\leq W_0-\kappa_W\tau_k-\delta k$, con $\tau_k$ el tiempo de movimiento
acumulado hasta el inicio de ese tubo, entonces $\dot W\leq-\kappa_W$ mientras
haya movimiento da $W\leq W_0-\kappa_W\tau_{k+1}-\delta k$ al final del tubo; y
si allí se admite una sustitución, la regla de admisión resta otros $\delta$,
de modo que $W\leq W_0-\kappa_W\tau_{k+1}-\delta(k+1)$ al inicio del tubo
$k+1$, que por composición de tubos también porta testigo seguro.

Al terminar la ejecución, $0\leq W\leq W_0-\kappa_WT_{\mathrm{mov}}-\delta N_{\mathrm{rev}}$,
y reordenando, $\kappa_WT_{\mathrm{mov}}+\delta N_{\mathrm{rev}}\leq W_0$. Con
$\kappa_W,\delta>0$ se obtienen las dos cotas del enunciado.

Lo demostrado es terminación, no vivacidad, y sobre el tiempo de movimiento,
no sobre la duración de la misión: nada en este argumento obliga a que el
objetivo se haya alcanzado al agotarse el presupuesto, ni acota cuánto puede
esperar la ejecución sin moverse.
\end{proof}

La existencia del testigo inicial no es gratuita. Si solo hay una parada
segura, se garantiza seguridad durante su dominio, no entrega. La regla que
compara \emph{upper bound} nuevo con presupuesto incumbente produce descenso
de presupuesto; solo una comparación con un \emph{lower bound} válido permite
afirmar mejora de coste frente a la alternativa anterior. Para márgenes
físicos $m_j\geq0$ con $\dot m_j\geq-\nu_j$, una condición local suficiente
de recuperación es
\[
 T_{\mathrm{rec}}=T_{\det}+T_{\mathrm{comm}}+T_{\mathrm{solve}}+T_{\mathrm{approach}}+T_{\mathrm{acquire}}+T_{\mathrm{commit}}
 <\min_{j:\nu_j>0}\frac{m_j(0)}{\nu_j}.
\]
No se suman indiscriminadamente márgenes de metros, newtons y julios. Los
fallos que destruyen todos los testigos quedan fuera de la hipótesis.

% La capa epistemologica delta se retira del documento principal: no fue
% implementada ni evaluada en el demostrador, y sus cinco ecuaciones numeradas
% no se citaban por su numero, lo que incumple la norma VIU. Su especificacion
% completa se conserva en supplementary/sections/epistemic-delta-gossip.tex.
\subsection{Extension no ensayada: actualizacion de informacion por diferencias}
\label{app:megajuego-epistemic-detail}

La informacion $\mathcal I$ de la Ecuacion~\eqref{eq:megajuego-state} admite una
extension en la que cada AMR transmite al vecino solo los hechos nuevos respecto
de la ultima version que ese vecino ha confirmado, en lugar del estado completo.
No se implemento ni se evaluo en el demostrador de esta memoria, de modo que no
sostiene ninguna afirmacion activa: no acredita menor trafico, decision mas
temprana ni robustez frente a perdida o retardo. Su especificacion, la regla de
admision y el diseno experimental necesario para contrastarla se conservan en el
material suplementario.

\subsection{Alcance y aporte}
\label{app:megajuego-scope}

La revisión externa identifica antecedentes directos de las piezas
genéricas: MPC distribuido mediante Smith \parencite*{barreiro2017distributed},
dinámicas poblacionales distribuidas, contacto cooperativo distribuido
mediante ADMM/DisCo, navegación con homotopías y juegos potenciales, VI
\emph{anytime} seguras, BeBOT y temporalización por alcanzabilidad. No se
reclama originalidad de esos bloques. La contribución candidata más coherente
es la contabilidad certificada de valor sobre revisiones atómicas físicamente
ejecutables: qué transición existe, qué estado y recursos consume, qué coste
común mejora, qué cota lo respalda y cómo conservar seguridad y progreso
mientras se calcula. Las pruebas de este documento demuestran dominios
concretos de esa arquitectura, no prioridad mundial sobre esos bloques
genéricos.

Para avanzar al siguiente nivel se necesita una campaña congelada con la
\emph{misma} planta de ruedas y contacto y el \emph{mismo} objetivo del juego
de integración y un control central de referencia, variando número de
agentes, contacto activo, pérdida de miembro, homotopía, batería, fricción,
retardo y presupuesto de cálculo; no se escogerán solo los casos que ambos
métodos completan como métrica principal. Esa campaña queda fuera del alcance
de esta memoria.
